% 00-sazetak.tex — Abstract (TASK §F.1). Numbers: data/raw/lab_session_01092026_020000/telemetry/run_log.csv.
Zadatak rada bio je preslikati zakret ljudskog zgloba na zglob robota Universal Robots u stvarnom vremenu uz jasno određene granice sigurnog ponašanja. Rad opisuje sustav koji to ostvaruje na kolaborativnom robotu UR3e i proširuje na vođenje cijele ruke. Sustav OptiTrack preko reflektirajućih markera prati nadlakticu, podlakticu i šaku kao tri kruta tijela, uz pet markera postavljenih na svakom segmentu. Iz njihovih se poza računa šest zglobnih kanala, a svaki kanal izravno vodi po jedan zglob robota bez inverzne kinematike. Traženo preslikavanje jednog anatomski podudarnog zgloba ostvareno je kanalom lakta na lakat robota (zglob J3). Između izračuna i robota nalaze se kauzalni One-Euro filtar i četveroslojni sigurnosni mehanizam s ograničenjem brzine zgloba, ograničenjem radnog raspona, zadržavanjem izlaza pri prekidu signala i zaustavljanjem.

Sustav je eksperimentalno vrednovan u laboratoriju kroz 21 pokus na stvarnom robotu. Proračun upravljačkog signala unutar petlje traje 3 do 4 ms na taktu od 125 Hz, ispod jednog postotka ukupnog dinamičkog odziva, pa računalni lanac nije usko grlo. Vjernost oponašanja pokreta nakon uklanjanja faznog pomaka iznosi 0,4 do 3,0$^\circ$ pri prirodnom tempu, odnosno do 6,2$^\circ$ pri izoliranom gibanju pojedinog zgloba, uz korelaciju 0,68 do 0,95. Ukupno dinamičko kašnjenje od pokreta ruke do odziva robota iznosi 385 ms na laktu. Razlaganjem po lancu na fazni pomak filtra (200 ms), sigurnosno ograničenje brzine (73 ms) i dinamiku servoa robota (112 ms) pokazuje se da odziv određuje filtriranje šuma, a ne mreža ili sam robot. Ispitana je i otpornost sustava na dodano mrežno kašnjenje, gubitak paketa, okluziju markera i djelovanje sigurnosnih slojeva, uz analizu granica upravljivosti u fizičkoj interakciji čovjeka i robota.
